%%%%%%%%%%%%%%%%%%%%%%%%%%%%%%%%%%%%%%%%%%%%%%%%%%%%%%%%%%%%%%%%%%%%%%
% Lớp 12 | Nguyên hàm, tích phân | Tích phân
% Dạng 2D4?2-4: Tích phân cơ bản hàm mũ, luỹ thừa
% Số câu: 21  (sinh tự động bởi tach_cau.py; từ nay thầy sửa trực tiếp file này)
%%%%%%%%%%%%%%%%%%%%%%%%%%%%%%%%%%%%%%%%%%%%%%%%%%%%%%%%%%%%%%%%%%%%%%

%%%%%%%%%%%%%%%% TRẮC NGHIỆM 4 PHƯƠNG ÁN %%%%%%%%%%%%%%%%

% nguồn: Phan1_Lop_12_Dai_so_Chuong_4
\begin{ex}%[2D4N2-4]%[ma:9144c11d]
	Tính $I=\displaystyle\int_0^1  3^x \mathrm{\,d}x$
	\choice 
		 {$I=2$}
		{\True$I=\dfrac{2}{\ln3}$} 
		{$I=\dfrac{1}{4}$} 
		 {$I=\dfrac{3}{\ln 3}$}
	\loigiai{ 
		$I=\displaystyle\int_0^1 3^x \mathrm{\,d}x=\dfrac{3^x}{\ln  3}\bigg|_0^1=\dfrac{3}{\ln3}-\dfrac{1}{\ln3}=\dfrac{2}{\ln3}$
	} 
\end{ex}

% nguồn: Phan1_Lop_12_Dai_so_Chuong_4
\begin{ex}%[2D4N2-4]%[ma:0564a2df]
	Tính $\displaystyle\int_1^4\left(\sqrt{x}+2^x \cdot \ln 2\right)  \mathrm{\,d}x$.
	 \choice
		 {$\dfrac{28}{3}$}
		 {\True $\dfrac{56}{3}$}
		{$\dfrac{58}{3}$} 
		{$\dfrac{29}{3}$} 
	\loigiai{ 
		Ta có $\displaystyle\int_1^4\left(\sqrt{x}+2^x \cdot \ln 2\right) \mathrm{\,d}x=\left.\left(\dfrac{2}{3} x \sqrt{x}+2^x\right)\right|_1 ^4=\left(\dfrac{2}{3} \cdot 4\cdot \sqrt{4}+2^4\right)-\left(\dfrac{2}{3} \cdot 1\cdot \sqrt{1}+2^1\right)=\dfrac{56}{3}$. 
	} 
\end{ex}

% nguồn: Phan1_Lop_12_Dai_so_Chuong_4
\begin{ex}%[2D4N2-4]%[ma:23ce7860]
	Tích phân $\displaystyle\int_0^3\mathrm{e}^x\mathrm{\,d}x$ có giá  trị bằng
	\choice 
		 {$1-\mathrm{e}^3$}
		 {\True $\mathrm{e}^3-1$}
		 {$\mathrm{e}^3+1$}
		{$\dfrac{477}{25}$} 
	\loigiai{Ta  có $\displaystyle\int_0^3\mathrm{e}^x\mathrm{\,d}x=\mathrm{e}^x\bigg|_0^3=\mathrm{e}^3-1$.}
\end{ex}

% nguồn: Phan1_Lop_12_Dai_so_Chuong_4
\begin{ex}%[2D4N2-4]%[ma:c4254916]
	Tìm họ nguyên hàm của  hàm số $f(x)=5^x$.
	\choice 
		{\True  $\displaystyle \int {f(x)}\mathrm{\,d}x=\dfrac{5^x}{\ln 5}+C$}
		 {$\displaystyle \int {f(x)}\mathrm{\,d}x=\dfrac{5^{x+1}}{x+1}+C$}
		{$\displaystyle  \int {f(x)}\mathrm{\,d}x=5^x+C$}
		{$\displaystyle \int {f(x)}\mathrm{\,d}x=5^x\ln 5+C$} 
	\loigiai{Ta có $\displaystyle \int {f(x)}\mathrm{\,d}x=\displaystyle \int  5^x\mathrm{\,d}x=\dfrac{5^x}{\ln 5}+C$.}
\end{ex}

% nguồn: Phan1_Lop_12_Dai_so_Chuong_4
\begin{ex}%[2D4N2-4]%[ma:fc526d5b]
	Cho  $\displaystyle\int_1^2 f(x) \mathrm{\,d}x=3$, $\displaystyle\int_2^3 f(x) \mathrm{\,d}x=-1$. Tính $\displaystyle\int_1^3 f(x) \mathrm{\,d}x$.
	\choice 
		{$-2$} 
		{$-4$} 
		 {\True$2$}
		{$4$} 
	\loigiai{ 
		Ta có  $\displaystyle\int_1^3 f(x) \mathrm{\,d}x=\displaystyle\int_1^2 f(x) \mathrm{\,d}x+\displaystyle\int_2^3 f(x) \mathrm{\,d}x=3+(-1)=2$.
	} 
\end{ex}

% nguồn: Phan1_Lop_12_Dai_so_Chuong_4
\begin{ex}%[2D4H2-4]%[ma:1c9cc6da]
	Nếu $\displaystyle\int_0^{\ln 3}{\left[f(x)+{\mathrm{e}^x}\right]\mathrm{\,d}x}=6$ thì  $\displaystyle\int_0^{\ln 3}{f(x)\mathrm{\,d}x}$ bằng
	 \choice
		{$6+\ln 3$} 
		{$6-\ln 3$} 
		{\True  $4$}
		{$8$} 
	 \loigiai{
		Ta có  \[\displaystyle\int_0^{\ln 3}{\left[f(x)+{\mathrm{e}^x}\right]\mathrm{\,d}x}=\displaystyle\int_0^{\ln 3}{f(x)\mathrm{\,d}x}+\displaystyle\int_0^{\ln 3}{{\mathrm{e}^x}\mathrm{\,d}x}=\displaystyle\int_0^{\ln 3}{f(x)\mathrm{\,d}x}+2.\]
		 Suy ra \[\displaystyle\int_0^{\ln 3}{f(x)\mathrm{\,d}x}=6-2=4.\]
	} 
\end{ex}

% nguồn: Phan1_Lop_12_Dai_so_Chuong_4
\begin{ex}%[2D4H2-4]%[ma:71e2f3b4]
	Cho  $\displaystyle\int_0^{\ln 2}{\big[2f(x)+{\mathrm{e}^x}\big]\mathrm{\,d}x}=5$. Khi đó $\displaystyle\int_0^{\ln 2}{f(x)\mathrm{\,d}x}$ bằng
	\choice 
		{$3$} 
		{$1$} 
		{\True $2$} 
		{$\dfrac52$} 
	\loigiai{ 
		Ta  có
		\allowdisplaybreaks 
		\begin{eqnarray*}
			\displaystyle\int_0^{\ln 2}{\big[2f(x)+{\mathrm{e}^x}\big]\mathrm{\,d}x}&=&2\displaystyle\int_0^{\ln 2}{f(x)\mathrm{\,d}x}+\displaystyle\int_0^{\ln  2}{{\mathrm{e}^x}\mathrm{\,d}x}\\
			&=&2\displaystyle\int_0^{\ln 2}{f(x)\mathrm{\,d}x}+ {\mathrm{e}^x} \Bigg|_0^{\ln 2}\\ 
			&=&2\displaystyle\int_0^{\ln  2}{f(x)\mathrm{\,d}x}+1.
		\end{eqnarray*}
		Mà $\displaystyle\int_0^{\ln  2}{\big[2f(x)+{\mathrm{e}^x}\big]\mathrm{\,d}x}=5\Rightarrow\displaystyle\int_0^{\ln 2}{f(x)\mathrm{\,d}x}=2$.
	 }
\end{ex}

% nguồn: Phan1_Lop_12_Dai_so_Chuong_4
\begin{ex}%[2D4H2-4]%[ma:ce471c6e]
	Cho hàm  số $ y=f(x)$ xác định và liên tục trên $\mathbb{R}$, thỏa mãn $\displaystyle\int_0^{\pi }{\left[f(x)+\sin x\right]\mathrm{\,d}x}=10$. Tính $I=\displaystyle\int_0^{\pi }{f(x)\mathrm{\,d}x}$.
	\choice 
		{$I=4$} 
		 {\True $I=8$}
		 {$I=12$}
		{$I=6$} 
	 \loigiai{
		 Ta có
		\allowdisplaybreaks 
		\begin{eqnarray*}
			\displaystyle\int_0^{\pi }{\left[f(x)+\sin x\right]\mathrm{\,d}x}&=&\displaystyle\int_0^{\pi }{f(x)\mathrm{\,d}x}+\displaystyle\int_0^{\pi }{\sin x\mathrm{\,d}x}\\ 
			 &=&\displaystyle\int_0^{\pi }{f(x)\mathrm{\,d}x}-\cos x\Bigg|_0^{\pi }\\
			&=&10. 
		\end{eqnarray*}
		Suy ra $\displaystyle\int_0^{\pi  }{f(x)\mathrm{\,d}x}=10+\left(\cos\pi-\cos 0\right)=8$.
	 }
\end{ex}

% nguồn: Phan1_Lop_12_Dai_so_Chuong_4
\begin{ex}%[2D4H2-4]%[ma:128483f7]
	Kết  quả tích phân $I=\displaystyle\int_0^1{5^x}\mathrm{\,d}x$ bằng
	 \choice
		 {\True $I=\dfrac4{\ln 5}$}
		{$I=4\ln  5$}
		{$I=5\ln 5$} 
		 {$I=\dfrac5{\ln 5}$}
	 \loigiai{
		Ta có \[I=\displaystyle\int_0^1{5^x}\mathrm{\,d}x=\dfrac{5^x}{\ln 5}  \Bigg|_0^1=\dfrac{5^1}{\ln 5}-\dfrac{5^0}{\ln 5}=\dfrac4{\ln 5}.\]
	} 
\end{ex}

% nguồn: Phan1_Lop_12_Dai_so_Chuong_4
\begin{ex}%[2D4V2-4]%[ma:c3ab3c23]
	Cho hàm  số $f(x)=\heva{&\mathrm{e}^x+m&&\text{ khi }x\geq 0\\&x^2\left(x^3+1\right)^3&&\text{ khi }x<0}$ (với $m$ tham số). Biết hàm số $f(x)$ liên tục trên $\mathbb{R}$ và tích phân $\displaystyle\int_{-1}^1f(x)\mathrm{\,d}x=a\cdot \mathrm{e}-\dfrac{b}{c}$ với $a,b,c\in\mathbb{Z}^{*}$; $\dfrac{a}{b}$ tối giản ($\mathrm{e}=2{,}718281\ldots$). Biểu thức $a+b+c+m$ có giá bằng
	 \choice
		 {$-11$}
		{\True $35$} 
		 {$13$}
		 {$36$}
	\loigiai{ 
		Vì  hàm số liên tục trên $\mathbb{R}$ nên hàm số liên tục tại điểm $x=0$.\\\
		Khi  đó
		 \allowdisplaybreaks
		\begin{eqnarray*}
			\lim\limits_{x\to 0^{+}}f(x)=\lim\limits_{x\to 0^{-}}f(x)=f(x)&\Leftrightarrow& \lim\limits_{x\to 0^{+}}\left(\mathrm{e}^x+m\right)=\lim\limits_{x\to 0^{+}}\big[x^2\left(x^3+1\right)^3\big]=1+m\\ 
			&\Leftrightarrow&  1+m=0\Leftrightarrow m=-1.
		\end{eqnarray*}
		Ta  có \[\displaystyle\int_{-1}^1f(x)\mathrm{\,d}x=\int_{-1}^0f(x)\mathrm{\,d}x+\int_0^1f(x)\mathrm{\,d}x=\int_0^1x^2\left(x^3+1\right)^3\mathrm{\,d}x+\int_0^1\left(\mathrm{e}^x-1\right)\mathrm{\,d}x=\mathrm{e}-\dfrac{23}{12}.\]
		Vậy $a=1$,  $b=23$, $c=12$ và $a+b+c+m=35$.
	} 
\end{ex}

% nguồn: Phan1_Lop_12_Dai_so_Chuong_4
\begin{ex}%[2D4V2-4]%[ma:b1e66c89]
	Cho hàm số $f(x)$ liên tục trên  $\mathbb{R}$ thỏa mãn $f(x)=\heva{&x+m,x\geq 0\\
	&\mathrm{e}^{2x},x<0}$ ($m$ là hằng số). Biết $\displaystyle\int_{-1}^2f(x)\mathrm{\,d}x=a+\dfrac{b}{\mathrm{e}^2}$ trong đó $a,b$ là các số  hữu tỉ. Tính $2a+4b^2$.
	\choice 
		 {\True $10$}
		 {$4$}
		{$80$} 
		{$40$} 
	 \loigiai{
		Do hàm số liên tục trên $\mathbb{R}$ nên  hàm số liên tục tại $x=0$.\\
		Khi  đó
		\[\lim\limits_{x\to 0^{-}}f(x)=\lim\limits_{x\to 0^{+}}f(x)=f(0)\Leftrightarrow  m=1.\]
		Ta có 
		 \allowdisplaybreaks
		\begin{eqnarray*}
			 \int_{-1}^2f(x)\mathrm{\,d}x&=&\int_{-1}^0f(x)\mathrm{\,d}x+\int_0^2f(x)\mathrm{\,d}x\\
			 &=&\int_{-1}^0\mathrm{e}^{2x}\mathrm{\,d}x+\int_0^2(x+1)\mathrm{\,d}x\\
			 &=&\dfrac{\mathrm{e}^{2x}}{2}\bigg|_{-1}^0+\left(\dfrac{x^2}{2}+x\right)\bigg|_0^2\\
			 &=&\dfrac{1}{2}-\dfrac{\mathrm{e}^{-2}}{2}+4\\
			 &=&\dfrac{9}{2}-\dfrac{1}{2\mathrm{e}^2}.
		\end{eqnarray*}
		 Do đó $a=\dfrac{9}{2};b=-\dfrac{1}{2}$.\\
		Vậy $2a+4b^2=10$. 
	 }
\end{ex}

% nguồn: Phan1_Lop_12_Dai_so_Chuong_4
\begin{ex}%[2D4V2-4]%[ma:9cc95c23]
	 Tính $\displaystyle\int_0^2\sqrt{x^2-2x+1} \mathrm{\,d} x$.
	\choice 
		 {$\dfrac{1}{2}$}
		 {$2$}
		{$\dfrac{5}{2}$} 
		{\True  $1$}
	 \loigiai{
		Ta có \begin{eqnarray*}
		\int_0^2\sqrt{x^2-2x+1} \mathrm{\,d}x&=&\int_0^2|x-1|\mathrm{\,d}x=\int_0^1(1-x) \mathrm{\,d}x+\int_1^2(x-1) \mathrm{\,d}x\\ 
		&=&\left.\left(x-\dfrac{x^2}{2}\right)\right|_0 ^1+\left.\left(\dfrac{x^2}{2}-x\right)\right|_1 ^2=1 
	\end{eqnarray*}
} 
\end{ex}


%%%%%%%%%%%%%%%% ĐÚNG/SAI %%%%%%%%%%%%%%%%

% nguồn: Phan1_Lop_12_Dai_so_Chuong_4
\begin{ex}%[2D4H2-4]%[ma:3f951030]
	Cho  hàm số $f(x)=\mathrm{e}^x$. Khi đó
	\choiceTF 
		{$\displaystyle\int_0^1  f(x)\mathrm{\,d}x=\mathrm{e}$}
		{$\displaystyle\int_0^1 f(2x+1)\mathrm{\,d}x=\mathrm{e}^3-\mathrm{e}$} 
		{\True  $\displaystyle\int_0^1 \big[f(x)-4x\big]\mathrm{\,d}x=\mathrm{e}-3$}
		{\True $\displaystyle\int_0^1  (x-1)f(x)\mathrm{\,d}x=2-\mathrm{e}$}
	 \loigiai{
		\begin{itemchoice}
			\itemch Ta  có $\displaystyle\int_0^1 f(x)\mathrm{\,d}x=\displaystyle\int_0^1\mathrm{e}^x\mathrm{\,d}x=\mathrm{e}-1$.
			\itemch  Ta có $\displaystyle\int_0^1 \mathrm{e}^{2x+1}\mathrm{\,d}x=\dfrac{1}{2}\mathrm{e}^{2x+1}\Big|_0^1=\dfrac{\mathrm{e}\left( \mathrm{e}^2-1 \right)}{2}$.
			\itemch Ta có $\displaystyle\int_0^1 \big[f(x)-4x\big]\mathrm{\,d}x=\displaystyle\int_0^1 \left( \mathrm{e}^x-4x \right)\mathrm{\,d}x=\mathrm{e}-3$. 
			\itemch  Ta có $\displaystyle\int_0^1 (x-1)f(x)\mathrm{\,d}x=\displaystyle\int_0^1 (x-1)\mathrm{e}^x\mathrm{\,d}x=2-\mathrm{e}$.
		\end{itemchoice}
	 }
\end{ex}

% nguồn: Phan1_Lop_12_Dai_so_Chuong_4
\begin{ex}%[2D4V2-4]%[ma:ff96c76e]
	Cho hàm  số $f(x)$ liên tục trên $\mathbb{R}$ thỏa mãn $f'(x)=x+\mathrm{e}^x$, $\forall x\in\mathbb{R}$.
	\choiceTF 
		 {$\displaystyle\int_1^2 f'(x)\mathrm{\,d}x=\mathrm{e}^2+\mathrm{e}+\dfrac{5}{2}$}
		{Diện tích hình phẳng giới hạn  bởi đồ thị của các đường $y=f'(x)$, $y=x+1$ và $x=2$ là $\dfrac{57}{13}$}
		{$f(x)=\mathrm{e}^x+x^2+C$} 
		{\True Khi $f(0)=4$ thì $\displaystyle\int_0^1 f(x)\mathrm{\,d}x=\dfrac{6\mathrm{e}+13}{6}$} 
	\loigiai{ 
	\begin{itemchoice}
		\itemch  Ta có \[\displaystyle\int_1^2 f'(x)\mathrm{\,d}x=\displaystyle\int_1^2 \left(x+\mathrm{e}^x\right) \mathrm{d}x= \left(\mathrm{e}^x+\dfrac{1}{2}x^2\right)\Bigg|^2_1=\mathrm{e}^2-\mathrm{e}+\dfrac{3}{2}.\]
		\itemch Phương  trình hoành độ giao điểm của các đường $y=f'(x)$, $y=x+1$ là
		\[x+\mathrm{e}^x=x+1 \Leftrightarrow \mathrm{e}^x=1 \Leftrightarrow x=0.\] 
		Trên đoạn $[0;2]$ ta có  $f'(x)-(x+1)=\mathrm{e}^x-1>0$. \\
		Diện tích hình phẳng  giới hạn là \[S=\displaystyle\int_0^2 \left(\mathrm{e}^x-1\right)\mathrm{\,d}x=\mathrm{e}^2-3.\]
		 \itemch $f(x)=\displaystyle\int f'(x)\mathrm{\,d}x=\mathrm{e}^x+\dfrac{1}{2}x^2+C$.
		\itemch Có  $f(x)=\displaystyle\int f'(x)\mathrm{\,d}x=\mathrm{e}^x+\dfrac{1}{2}x^2+C$. \\
		Với $f(0)=4 \Leftrightarrow \mathrm{e}^0+\dfrac{1}{2}\cdot0^2+C=4 \Leftrightarrow  C=4$. \\
		Suy ra $f(x)=\mathrm{e}^x+\dfrac{1}{2}x^2+4$. \\ 
		Khi đó $\displaystyle\int_0^1 f(x)\mathrm{\,d}x=\mathrm{e}+\dfrac{6}{13}=\dfrac{6\mathrm{e}+13}{6}$. 
	\end{itemchoice}
	} 
\end{ex}


%%%%%%%%%%%%%%%% TRẢ LỜI NGẮN %%%%%%%%%%%%%%%%

% nguồn: Phan2_Lop_12_Dai_so_Chuong_4
\begin{bt}%[2D4N2-4]%[ma:1974ac27]
	 [Dự án đề kiểm tra Toán 12 HKII NH24-25- Hector Tran]
	Cho $\displaystyle\int_{0}^{2} e^{4x} \mathrm{\,d}x = \dfrac{1}{4} e^m  + n$ với $m$, $n \in \mathbb{R}$. Tính $m+16n$?
	\shortans{4} 
	\loigiai{Ta có $\displaystyle\int_{0}^{2} e^{4x} \mathrm{\,d}x = \frac{1}{4}e^{4x} \Big|_{0}^{2}= \frac{1}{4}e^8 - \frac{1}{4}= \frac{1}{4}  e^m + n$.\\
		 Ta được $m=8$ và $n = -\frac{1}{4}$.\\
	Vậy $m+16n =  8 + 16\left(-\frac{1}{4}\right) = 8 - 4 = 4$.\\}
\end{bt}

% nguồn: Phan2_Lop_12_Dai_so_Chuong_4
\begin{bt}%[2D4H2-4]%[ma:72024449]
	Giả sử $\forall a, b \in \mathbb{R}, a<0<b, \displaystyle\int_a^b|x|^7\, \mathrm{\,d}x=ma^8+n b^8$ trong đó $m, n$ là các hằng số  thực (không phụ thuộc vào $a$ và $b$). Giá trị của biểu thức $P=m-5 n$ là bao nhiêu?
	\shortans{-0,5} 
	 \loigiai{
		Với mọi  $a, b \in \mathbb{R}, a<0<b$, ta có
		\allowdisplaybreaks 
		\begin{eqnarray*}
			\displaystyle\int_a^b|x|^7\mathrm{\,d}x  &=&\displaystyle\int_a^0|x|^7\mathrm{\,d}x+\displaystyle\int_0^b|x|^7\mathrm{\,d}x \\
			&=&\displaystyle\int_a^0\left(-x^7\right)\mathrm{\,d}x+\displaystyle\int_0^b x^7\mathrm{\,d}x\\ 
			&=&\dfrac{1}{8} a^8+\dfrac{1}{8} b^8. 
		\end{eqnarray*}
		 Do đó $P=\dfrac{1}{8}-\dfrac{5}{8}=-0{,}5$.
	 }
\end{bt}

% nguồn: Phan2_Lop_12_Dai_so_Chuong_4
\begin{bt}%[2D4H2-4]%[ma:559ec46b]
	$\displaystyle\int_0^1 \dfrac{3^{x-2}}{2^{2  x}} \mathrm{\,d}x$ có giá trị bằng bao nhiêu? (viết kết quả dưới dạng số thập phân và làm tròn đến hàng phần mười).
	\shortans{0,1} 
	\loigiai{ 
	Ta có $\displaystyle\int_0^1 \dfrac{3^{x-2}}{2^{2 x}} \mathrm{\,d} x=\dfrac{1}{9} \displaystyle\int_0^1 \left(\dfrac{3}{4}\right)^x  \mathrm{\,d} x=\left.\dfrac{\left(\dfrac{3}{4}\right)^x}{9 \ln\dfrac{3}{4}}\right|_0 ^1=\dfrac{3}{36 \ln\dfrac{3}{4}}-\dfrac{1}{9 \ln\dfrac{3}{4}} =-\dfrac{1}{36 \ln\dfrac{3}{4}}\approx 0{,}1$.}
\end{bt}


%%%%%%%%%%%%%%%% TỰ LUẬN %%%%%%%%%%%%%%%%

% nguồn: Phan2_Lop_12_Dai_so_Chuong_4
\begin{bt}%[2D4H2-4]%[ma:2c781ff9]
	Thực hiện các yêu cầu sau  đây:
	\begin{enumerate}
		\item Cho $F(x)$ là một nguyên hàm của hàm số $  f(x)={3^x}$ trên $\mathbb{R}$ thỏa mãn $F(0)=0$. Tính $F(1)$.
		\item Cho $F(x)$  là một nguyên hàm của hàm $ f(x)={2^x}$ và $F(0)=0$. Tính $F(1)$.
	\end{enumerate}
	 \loigiai{
		\begin{enumerate}
			\item  Ta có \[F(1)-F(0)=\displaystyle\int_0^1{f(x)\mathrm{\,d}x}=\displaystyle\int_0^1{3^x\mathrm{\,d}x=\dfrac{3^x}\ln 3}\Bigg|_0^1=\dfrac3{\ln 3}-\dfrac1{\ln 3}=\dfrac2{\ln 3}.\]
			 Suy ra \[F(1)=\dfrac2{\ln 3}+F(0)=\dfrac2{\ln 3}+0=\dfrac2{\ln 3}.\]
			 \item Ta có \[F(1)-F(0)=\displaystyle\int_0^1{f(x)\mathrm{\,d}x}=\displaystyle\int_0^1{{2^x} \, \mathrm{\,d}x}=\dfrac{2^x}{\ln2} \Bigg|_0^1=\dfrac{{2^1}-{2^0}}{\ln2}=\dfrac1{\ln2}= \log_2\mathrm e.\]
			Vậy  $F(1)= \log_2\mathrm e$.
		\end{enumerate}
	} 
\end{bt}

% nguồn: Phan2_Lop_12_Dai_so_Chuong_4
\begin{bt}%[2D4H2-4]%[ma:c42ddf31]
	Tính các tích phân sau: 
	\begin{listEX}[2]
		\item[a)]  $\displaystyle \int_{0}^{1}(4x^3+x)\mathrm{\,d}x$;
		\item[b)] $\displaystyle  \int_{1}^{2}\dfrac{x-2}{x^2}\mathrm{\,d}x$;
		 \item[c)] $\displaystyle \int_{0}^{4}2^{2x}\mathrm{\,d}x$;
		\item[d)]  $\displaystyle \int_{0}^{1}(\mathrm{e}^{x-1}+2^{x+1})\mathrm{\,d}x$.
	\end{listEX}
	 \loigiai{
		\begin{enumerate}
			 \item $\displaystyle \int_{0}^{1}(4x^3+x)\mathrm{\,d}x=\left(x^4+\dfrac{x^2}{2}\right)\bigg|_{0}^{1}=\dfrac{3}{2}$;
			\item $\displaystyle  \int_{1}^{2}\dfrac{x-2}{x^2}\mathrm{\,d}x=\displaystyle \int_{1}^{2}\dfrac{1}{x}-\dfrac{2}{x^2}\mathrm{\,d}x=\left(\ln x+\dfrac{2}{x}\right)\bigg|_{1}^{2}=\ln2+1-2=\ln2-1$;
			 \item $\displaystyle \int_{0}^{4}2^{2x}\mathrm{\,d}x=\dfrac{2^{2x}}{2\ln2}\bigg|_{0}^{4}=\dfrac{2^8}{2\ln2}-\dfrac{1}{2\ln2}=\dfrac{255}{2\ln2}$;
			\item $\displaystyle \int_{0}^{1}(\mathrm{e}^{x-1}+2^{x+1})\mathrm{\,d}x=\left(\mathrm{e}^{x-1}+\dfrac{2^{x+1}}{\ln2}\right)\bigg|_{1}^{2}=\mathrm{e}+\dfrac{2^3}{\ln2}-(1+\dfrac{2^2}{\ln2})=\dfrac{4}{\ln2}+\mathrm{e}-1$. 
		\end{enumerate}
	 }
\end{bt}

% nguồn: Phan2_Lop_12_Dai_so_Chuong_4
\begin{bt}%[2D4H2-4]%[ma:29bd03be]
	Tính các tích phân  sau đây:
	\begin{enumEX}{3}
		 \item $\displaystyle\int_0^3{3x^2\mathrm{\,d}x}$
		\item $\displaystyle\int_0^{\tfrac{\pi }3}{\sin x\mathrm{\,d}x}$ 
		 \item $\displaystyle\int_1^{\mathrm{e}^2}{\dfrac1u} \mathrm{\,d}u$
		\item $\displaystyle\int_0^{\tfrac{\pi  }4}{\dfrac1{{{\cos }^2}x}\mathrm{\,d}x}$
		\item $\displaystyle\int_{-2}^4{\left(x+1\right)\left(x-1\right)}\mathrm{\,d}x$ 
		\item $\displaystyle\int_1^2{\dfrac{{x^2}-2x+1}x}  \mathrm{\,d}x$
		\item  $\displaystyle\int_0^{\tfrac\pi2}{\left(3\sin x-2\right)} \mathrm{\,d}x$
		 \item $\displaystyle\int_0^{\tfrac\pi2}{\dfrac{{{\sin }^2}x}{1+\cos x}} \mathrm{\,d}x$
		\item $\displaystyle\int_1^4{\left(\sqrt  x+2\right)\left(\sqrt x-2\right)}\mathrm{\,d}x$
		 \item $\displaystyle\int_2^3{\dfrac{{x^2}-3x+1}{x-1}\mathrm{\,d}x}$
		\item $\displaystyle\int_0^{\tfrac\pi2}{\left(2{{\sin  }^2}x+3\right)}\mathrm{\,d}x$
		\item $\displaystyle\int_0^{\tfrac{\pi }4}{\dfrac{\cos x+\cos3x}{\cos  x}}\mathrm{\,d}x$
		\item $\displaystyle\int_0^{\tfrac{\pi }4}{\dfrac{\cos2x}{\cos x\left(1+\tan x\right)}}\mathrm{\,d}x$ 
		 \item $\displaystyle\int_{\tfrac{\pi }6}^{\tfrac{\pi }4}{\left(\dfrac1{\cos^2x}-\dfrac1{\sin^2x}\right)\mathrm{\,d}x}$
		\item  $\displaystyle\int_0^1{\left({2^{2x}}\cdot3^{x-1}\right)\mathrm{\,d}x}$
	\end{enumEX}
	 \loigiai{
		\begin{enumerate}
			\item Ta  có \[\displaystyle\int_0^3{3x^2\mathrm{\,d}x}=x^3\Bigg|_0^3={3^3}-{0^3}=27.\]
			\item Ta có \[\displaystyle\int_0^{\tfrac{\pi }3}{\sin x\mathrm{\,d}x}=-\cos x\Bigg|_0^{\tfrac{\pi }3}=-\left(\cos\dfrac{\pi  }3-\cos0\right)=\dfrac12.\]
			 \item Ta có \[\displaystyle\int_1^{\mathrm{e}^2}{\dfrac1u} \mathrm{\,d}u=\ln|u |\Bigg|_1^{\mathrm{e}^2}=\ln {\mathrm{e}^2}-\ln 1=2.\]
			\item Ta có \[\displaystyle\int_0^{\tfrac{\pi }4}{\dfrac1{{{\cos  }^2}x}\mathrm{\,d}x}=\tan x \Bigg|_0^{\tfrac{\pi }4}=1.\]
			\item Ta  có \[\displaystyle\int_{-2}^4{\left(x+1\right)\left(x-1\right)}\mathrm{\,d}x=\displaystyle\int_{-2}^4{\left({x^2}-1\right)}\mathrm{\,d}x=\left(\dfrac{x^3}3-x\right) \Bigg|_{-2}^4=18.\]
			\item Ta có \[\displaystyle\int_1^2{\dfrac{{x^2}-2x+1}x} \mathrm{\,d}x=\displaystyle\int_1^2{\left(x-2+\dfrac1x\right)  }\mathrm{\,d}x=\left(\dfrac{x^2}2-2x+\ln x\right) \Bigg|_1^2=-\dfrac12+\ln 2.\]
			\item Ta có \[\displaystyle\int_0^{\tfrac\pi2}{\left(3\sin x-2\right)} \mathrm{\,d}x=\left(-3\cos x-2x\right) \Bigg|_0^{\tfrac\pi2}=-3\left(\cos\dfrac\pi2-\cos  0\right)-2\left(\dfrac\pi2-0\right)=3-\pi .\]
			\item Ta có \[\displaystyle\int_0^{\tfrac\pi2}{\dfrac{{{\sin }^2}x}{1+\cos x}} \mathrm{\,d}x=\displaystyle\int_0^{\tfrac\pi2}{\dfrac{1-{{\cos  }^2}x}{1+\cos x}} \mathrm{\,d}x=\displaystyle\int_0^{\tfrac\pi2}{\left(1-\cos x\right)} \mathrm{\,d}x=\left(x-\sin x\right) \Bigg|_0^{\tfrac\pi2}=\dfrac\pi2-1.\]
			\item Ta có \[\displaystyle\int_1^4{\left(\sqrt x+2\right)\left(\sqrt  x-2\right)}\mathrm{\,d}x=\displaystyle\int_1^4{\left(x-4\right)}\mathrm{\,d}x=\left(\dfrac{x^2}2-4x\right) \Bigg|_1^4=\dfrac{-9}2.\]
			\item Ta có \[\displaystyle\int_2^3{\dfrac{{x^2}-3x+1}{x-1}\mathrm{\,d}x}=\displaystyle\int_2^3{x-2-\dfrac1{x-1}\mathrm{\,d}x}=\left(\dfrac{x^2}2-2x-\ln|x-1  |\right)_2^3=\dfrac12-\ln 2.\]
			\item Ta có \[\displaystyle\int_0^{\tfrac\pi2}{\left(2{{\sin }^2}x+3\right)}\mathrm{\,d}x=\displaystyle\int_0^{\tfrac\pi2}{\left(1-\cos2x+3\right)}\mathrm{\,d}x=\displaystyle\int_0^{\tfrac\pi2}{\left(4-\cos2x\right)}\mathrm{\,d}x=\left(4x-\dfrac12\sin2x\right) \Bigg|_0^{\tfrac\pi2}=2\pi .\] 
			\item Ta có \[\displaystyle\int_0^{\tfrac{\pi }4}{\dfrac{\cos x+\cos3x}{\cos x}}\mathrm{\,d}x=\displaystyle\int_0^{\tfrac{\pi  }4}{\dfrac{2\cos x.\cos2x}{\cos x}}\mathrm{\,d}x=\displaystyle\int_0^{\tfrac{\pi }4}{2\cos2x}\mathrm{\,d}x=\left(\sin2x\right) \Bigg|_0^{\tfrac{\pi }4}=1.\]
			\item  Ta có
			\begin{align*}
				\displaystyle\int_0^{\tfrac{\pi }4}{\dfrac{\cos2x}{\cos x\left(1+\tan  x\right)}}\mathrm{\,d}x&=\displaystyle\int_0^{\tfrac{\pi }4}{\dfrac{\left(\cos x-\sin x\right)\left(\cos x+\sin x\right)}{\cos x+\sin x}}\mathrm{\,d}x\\
				&=\displaystyle\int_0^{\tfrac{\pi }4}{\left(\cos  x-\sin x\right)}\mathrm{\,d}x=\left(\sin x+\cos x\right) \Bigg|_0^{\tfrac{\pi }4}=\sqrt2-1.
			\end{align*}
			 \item Ta có \[\displaystyle\int_{\tfrac{\pi }6}^{\tfrac{\pi }4}{\left(\dfrac1{\cos^2x}-\dfrac1{\sin^2x}\right)\mathrm{\,d}x}=\left(\tan x+\cot x\right) \Bigg|_{\tfrac{\pi }6}^{\tfrac{\pi }4}=2-\dfrac{4\sqrt3}3.\]
			\item Ta có  \[\displaystyle\int_0^1{\left({2^{2x}}\cdot3^{x-1}\right)\mathrm{\,d}x}=\dfrac13\displaystyle\int_0^1{{{12}^x} \, \mathrm{\,d}x}=\dfrac13\cdot\dfrac{{{12}^x}}{\ln12}\Bigg|_0^1=\dfrac4{\ln12}-\dfrac1{3 \ln12}=\dfrac{11}{3 \ln12}.\]
		\end{enumerate}
	} 
\end{bt}

% nguồn: Phan2_Lop_12_Dai_so_Chuong_4
\begin{bt}%[2D4H2-4]%[ma:babeaa0d]
	Tính các tích  phân sau đây:
	\begin{enumEX}{3}
		\item $\displaystyle\int_0^1{\left(4x^3+3x^2-2\right)\mathrm{\,d}x}$ 
		 \item $\displaystyle\int_1^4{\dfrac1{2\sqrt x}\mathrm{\,d}x}$
		\item $\displaystyle\int_1^3{{x^{\sqrt2}}\mathrm{\,d}x}$ 
		 \item $\displaystyle\int_1^2\dfrac1{x^3}\mathrm{\,d}x$
		 \item $\displaystyle\int_1^3{{x^{\tfrac23}}}\mathrm{\,d}x$
		 \item $\displaystyle\int_1^8{\sqrt[3]x}\mathrm{\,d}x$
		\item $\displaystyle\int_{-\mathrm e}^{-1}{\dfrac2x\mathrm{\,d}x.}$ 
		 \item $\displaystyle\int_1^{\mathrm e}{\dfrac7{3x}\,\mathrm{\,d}x}$
		\item  $\displaystyle\int_0^1{\left(2\cdot3^x-5{\mathrm{e}^x}\right)\,\mathrm{\,d}x}$
	\end{enumEX}
	 \loigiai{
		\begin{enumerate}
			\item  Ta có
			\begin{align*}
				\displaystyle\int_0^1{\left(4x^3+3x^2-2\right)\mathrm{\,d}x}&=\displaystyle\int_0^1{4x^3\mathrm{\,d}x}+\displaystyle\int_0^1{3x^2\mathrm{\,d}x}-\displaystyle\int_0^1{2\mathrm{\,d}x}=x^4\Bigg|_0^1+{x^3}\Bigg|_0^1-2x\Bigg|_0^1\\ 
				&=\left({1^4}-{0^4}\right)+\left({1^3}-{0^3}\right)-\left(2\cdot1-2\cdot0\right)=0. 
			\end{align*}
			\item Ta có \[\displaystyle\int_1^4{\dfrac1{2\sqrt  x}\mathrm{\,d}x} =\dfrac12\displaystyle\int_1^4{{x^{-\tfrac12}}\mathrm{\,d}x}{x^{\tfrac12}}\Bigg|_1^4=\sqrt x\Bigg|_1^4=\sqrt4-\sqrt1=1.\]
			\item Ta  có \[\displaystyle\int_1^3{{x^{\sqrt2}}\mathrm{\,d}x}=\dfrac{{x^{\sqrt2+1}}}{\sqrt2+1}\Bigg|_1^3=\dfrac{{3^{\sqrt2+1}}-1}{\sqrt2+1}.\]
			\item Ta có \[\displaystyle\int_1^2\dfrac1{x^3}\mathrm{\,d}x=\displaystyle\int_1^2{{x^{-3}}}\mathrm{\,d}x=\dfrac{{x^{-2}}}{-2} \Bigg|_1^2=\dfrac{{2^{-2}}}{-2}-\dfrac{{1^{-2}}}{-2}=\dfrac38.\] 
			 \item Ta có \[\displaystyle\int_1^3{{x^{\tfrac23}}}\mathrm{\,d}x=\dfrac35{x^{\tfrac53}} \Bigg|_1^3=\dfrac35\cdot\left({3^{\tfrac53}}-1\right)=\dfrac35\cdot\left(\sqrt[3]{3^5}-1\right).\]
			\item Ta  có \[\displaystyle\int_1^8{\sqrt[3]x}\mathrm{\,d}x=\displaystyle\int_1^8{{x^{\tfrac13}}}\mathrm{\,d}x=\dfrac34\cdot{x^{\tfrac43}} \Bigg|_1^8=\dfrac34.\left({8^{\tfrac43}}-1\right)=\dfrac34\left(\sqrt[3]{8^4}-1\right)=\dfrac34\cdot\left(16-1\right)=\dfrac{45}4.\]
			 \item Ta có \[\displaystyle\int_{-\mathrm e}^{-1}{\dfrac2x\mathrm{\,d}x=2\ln|x |\Bigg|_{-\mathrm e}^{-1}}=2\ln|-1 |-2\ln|-\mathrm e|=-2.\]
			\item Ta có \[\displaystyle\int_1^{\mathrm e}{\dfrac7{3x}\,\mathrm{\,d}x}=\dfrac73\ln|x |\Bigg|_1^{\mathrm e}=\dfrac73\cdot\ln\mathrm  e-\dfrac73\cdot\ln 1=\dfrac73.\]
			\item Ta có \[\displaystyle\int_0^1{\left({{2\cdot3}^x}-5{\mathrm{e}^x}\right)\,\mathrm{\,d}x}=\left(\dfrac{2\cdot3^x}{\ln 3}-5{\mathrm{e}^x}\right) \Bigg|_0^1=\left(\dfrac6{\ln 3}-5\mathrm  e\right)-\left(\dfrac2{\ln 3}-5\right)=\dfrac4{\ln 3}-5\mathrm e+5.\]
		\end{enumerate}
	} 
\end{bt}
